\chapter{Hardware Concepts for Transactional Memory}
\index{memory hierarchy}\index{cache coherence}\index{MESI}\index{out-of-order execution}\index{speculation}\index{store buffer}\index{memory ordering}\index{TSO}\index{fence}\index{false conflict}\index{HTM}\index{capacity abort}\index{fallback path}\index{commit-publish}

Transactional memory runs on hardware that was not designed for it. The
machine presents a shared address space, but underneath there are private
caches, a coherence protocol, store buffers, and speculative execution, and
every one of those features shapes what a TM design can and cannot do. A
designer who treats memory as a single flat array will be surprised by false
conflicts, capacity aborts, and store-buffer reorderings the first time a real
system runs.

This chapter builds the hardware vocabulary used everywhere else in the book:
the memory hierarchy, cache coherence, out-of-order execution and speculation,
and the memory-ordering models that operating systems and compilers expose.
The point is not to design a processor but to know, for each later algorithm,
which hardware fact makes it correct and which one makes it cheap.

\section{Multicore Architectures and Memory Hierarchies}

The first fact a transactional memory designer must internalize is that modern
shared-memory machines are not single memories protected by a bus. They are
hierarchies of private and shared caches with a protocol that keeps those
caches approximately consistent. A multicore processor presents a shared
address space, but the physical reality is per-core L1 and L2 caches, a shared
L3, and a coherence protocol that moves cache lines between cores. The TM
substrate is this protocol.

Cache coherence is the mechanism that keeps copies of the same line
consistent. The most commonly taught protocol is the MESI family, named for
its four states: Modified, Exclusive, Shared, and Invalid. A line in the
Modified state is dirty and exists only in that cache. A line in Exclusive
state is clean and private. A line in Shared state is clean and may be
replicated. Invalid means the copy is not usable. From a transactional memory
viewpoint, the key transition is invalidation: when one core writes a line,
the coherence protocol invalidates that line in all other caches. An
invalidation that hits a line inside another transaction's read- or write-set
is a \emph{conflict signal}.

Coherence comes in two broad families. Snooping-based protocols broadcast
coherence traffic on a shared interconnect. Directory-based protocols keep a
directory entry per line and send targeted invalidations. Both deliver the
same logical guarantee: at any moment, either one writer has exclusive access,
or multiple readers share the line. The important caveat for later chapters is
that GPUs do not provide this coherent private cache model across streaming
multiprocessors. GPU global memory is coherent only at L2, and the private L1
caches are not kept coherent across SMs. That absence reshapes GPU TM design
in Chapter~17.

Conflict detection granularity is the cache line, not the word. Two
independent variables can share a line and therefore appear to conflict even
when no thread touches the same word. This false conflict is not an
implementation bug; it is the coherence protocol working exactly as designed.
Hardware TM inherits that granularity. Software TM can choose finer granularity,
but only by paying metadata and instrumentation costs that hardware gets for
free.

\section{Out-of-Order Execution and Speculation}

Modern cores do not execute instructions in program order. They fetch, decode,
and issue instructions into a reorder buffer, execute them when operands are
ready, and commit results in program order. Branch prediction permits the core
to execute speculatively past unresolved branches. If a branch prediction is
wrong, the core squashes the speculative work and restarts from the correct
path. This is ordinary processor speculation.

The reorder buffer holds register state, not memory state. A speculative store
does not become architecturally visible until the branch resolves and the
store retires. Hardware transactional memory extends this principle to a
region of instructions. A transaction's speculative loads and stores are
tracked in the cache, not in the ROB, and they become visible only at commit.
If the transaction aborts, the cache lines are invalidated or restored to
their pre-transaction values. The parallel between branch misprediction
recovery and transactional abort is exact: both discard speculative work and
restart from a stable point.

The capacity limit follows from this placement. The transaction working set
must fit in the private cache, typically L1 or L2, because the speculative
data has no other home. A transaction that touches more lines than the cache
can track is capacity-aborted. Large transactions therefore need either a
software TM that buffers state in memory, or a hardware fallback path that
takes a lock and executes the transaction non-speculatively.

Figure~\ref{fig:pipeline-squash} shows the lifecycle of a speculative store in
a simplified pipeline: the store is executed speculatively in the memory
stage, but if the branch mispredicts, the writeback is squashed and the store
never becomes visible.

\begin{figure}[htbp]
\centering
\begin{tikzpicture}[
font=\small,
state/.style={draw,rectangle,rounded corners,align=center,
minimum width=2.2cm,minimum height=0.9cm},
edge/.style={-{Stealth},thick}
]
\node[state] (fetch) {Fetch};
\node[state, right=0.85cm of fetch] (decode) {Decode};
\node[state, right=0.85cm of decode] (execute) {Execute};
\node[state, right=0.85cm of execute] (mem) {Memory\\speculative\\store};
\node[state, right=0.85cm of mem] (writeback) {Writeback};
\draw[edge] (fetch) -- (decode);
\draw[edge] (decode) -- (execute);
\draw[edge] (execute) -- (mem);
\draw[edge] (mem) -- (writeback);
% commit annotation routed above the boxes
\draw[edge, dashed] (mem.north) -- ++(0,0.6) -|
  node[pos=0.25, above, font=\footnotesize]
  {commit: store becomes visible} (writeback.north);
\draw[edge, dashed, red] (mem.south) -- ++(0,-0.8)
node[below, red, font=\footnotesize] {branch mispredict: squash};
\end{tikzpicture}
\caption{Speculative store in a simplified pipeline. The store reaches the
memory stage, but if a prior branch mispredicts, writeback is squashed and
the store never becomes architecturally visible. Hardware TM applies the
same idea to a multi-instruction region.}
\label{fig:pipeline-squash}
\end{figure}

\section{Store Buffers and Memory Ordering}

Between the core and the cache sits one more small, private structure: the
\emph{store buffer}. A store buffer is a small per-core FIFO that holds
completed stores until they are written to the memory hierarchy. Store
buffers improve performance by letting the core continue while stores drain
in the background. But a store buffer also relaxes write-to-read ordering,
and it is the reason that two cores can disagree about the order in which
two stores became visible.

A \emph{memory consistency model} is the contract that says which such
reorderings are legal. The contract is easiest to state with \emph{litmus
tests}: tiny two-thread programs with a question attached (``can this
outcome occur?''). Each litmus test exercises one pair of operations---store
then load, load then load---and the model answers yes or no for each. This
section walks through the models a TM author actually meets, from the
strongest (sequential consistency) through TSO, the model of x86-64
processors, to the weakly ordered ARM and POWER families, and names the
anomaly each model permits. The pattern to watch for is always the same: the
model permits an interleaving that sequential consistency forbids, that
forbidden interleaving breaks a transaction's atomicity, and a fence or a
stronger atomic closes the hole.

\subsection{Sequential Consistency}

Lamport's sequential consistency (SC) is the strongest shared-memory
contract \cite{lamport79sc}: the result of any execution is the same as if
the operations of all cores had been executed in \emph{some} sequential
order, and the operations of each individual core appear in that order in
program order. Figure~\ref{fig:sc-model} shows the shape of such an order
for a two-thread execution. There is nothing to add and nothing to take
away: every execution is one interleaving, and a transaction maps onto it
naturally as a contiguous segment.

Under SC no litmus test has a surprising outcome. The price is that the
hardware may not reuse a store buffer to hide write latency: a store must
reach the coherent memory order before the next load of the same core
completes. No commercial multicore processor is sequentially consistent,
precisely because that prohibition costs performance. SC is nevertheless
the model worth naming first, because it is the reference against which
every weaker model is defined---and because the correctness conditions of
Chapter~5 (serializability, opacity) are stated in its image.

\begin{figure}[htbp]
\centering
\begin{tikzpicture}[
font=\small,
op/.style={draw,rectangle,rounded corners,align=center,
minimum width=5.6cm,minimum height=0.7cm}
]
\node[op] (o1) {1. \; \texttt{W} $A{=}1$ \hfill (Core 0)};
\node[op, below=0.22cm of o1] (o2) {2. \; \texttt{W} $B{=}1$ \hfill (Core 0)};
\node[op, below=0.22cm of o2] (o3) {3. \; \texttt{R} $B \to 1$ \hfill (Core 1)};
\node[op, below=0.22cm of o3] (o4) {4. \; \texttt{R} $A \to 1$ \hfill (Core 1)};
\draw[dashed,rounded corners]
  ($(o1.north west)+(-0.35,0.35)$) rectangle ($(o4.south east)+(0.35,-0.3)$);
\node[font=\footnotesize, above=0.42cm of o1] {one global order (an interleaving)};
\end{tikzpicture}
\caption{Sequential consistency. All operations from all cores are woven
into a single global order; each core's operations keep their program
order inside it. No outcome of a litmus test can deviate from this shape,
so no anomaly exists by construction.}
\label{fig:sc-model}
\end{figure}

\subsection{Total Store Order (x86-64, SPARC)}
\label{sec:tso-model}

The memory model of x86-64 (and SPARC's TSO variant) can be described in
one sentence: \emph{the store buffer is the only source of reordering}
\cite{sewell10cacm,adve96tutorial}. Loads and stores keep their order with
respect to each other---load after load, store after store, load then
store---except for one pair: a \emph{store followed by a load to a
different address} may appear to complete out of order, because the load
can execute while the earlier store is still sitting in the store buffer.
A load to the \emph{same} address as a buffered store is forwarded from
the buffer itself (store-to-load forwarding), so the core always sees its
own writes in order. The model remains multi-copy atomic: a store becomes
visible to all other cores at once, in one global memory order.

Figure~\ref{fig:tso-model} shows the one anomaly this relaxation permits,
in its purest form. Both threads store to one variable and then load the
other. Each load bypasses its own store buffer, reads the shared memory,
and finds $0$. The outcome $r_1{=}r_2{=}0$---both threads believe they went
first---is forbidden under SC and allowed under TSO. This litmus test is
the classic \emph{store-buffering} (SB) test, also known as Dekker's
example, and closing it requires a full \texttt{mfence} between the store
and the load.

Two further consequences follow from the same relaxation. First, a core
observes its own stores before anyone else does, so a value read back
locally is not evidence that other cores can see it. Second, stores drain
from the buffer one at a time, so other cores can observe a multi-word
write in a partially drained state---the torn view that Section~\ref{sec:tso-example}
exploits. What TSO \emph{does} guarantee is store-store and load-load
order: the message-passing idiom of Section~\ref{sec:weak-model} is safe
on x86-64 with no hardware fence at all, which is why a protocol tuned on
x86 can silently hide a missing barrier until it is ported to ARM.

\begin{figure}[htbp]
\centering
\begin{tikzpicture}[
font=\small,
core/.style={draw,rectangle,rounded corners,align=center,
minimum width=4.4cm,minimum height=1.1cm},
buf/.style={draw,rectangle,align=center,
minimum width=4.4cm,minimum height=1.0cm},
mem/.style={draw,rectangle,rounded corners,align=center,
minimum width=11.2cm,minimum height=1.0cm},
edge/.style={-{Stealth},thick},
lab/.style={font=\footnotesize,fill=white,inner sep=1.5pt}
]
\node[core] (core0) at (0,0)
  {Core 0\\\texttt{W} $A{=}1$; \ then \texttt{R} $B \to r_1$};
\node[core] (core1) at (6.6,0)
  {Core 1\\\texttt{W} $B{=}1$; \ then \texttt{R} $A \to r_2$};
\node[buf]  (buf0)  at (0,-2.1) {store buffer\\$A{=}1$ pending};
\node[buf]  (buf1)  at (6.6,-2.1) {store buffer\\$B{=}1$ pending};
\node[mem]  (memory) at (3.3,-4.2)
  {shared memory: \ $A{=}0$, $B{=}0$ \ (nothing drained yet)};

\draw[edge] (core0) -- (buf0);
\draw[edge] (core1) -- (buf1);
\draw[edge,dashed] (buf0) -- (memory.north -| buf0)
  node[lab, midway, right=2pt] {drains later};
\draw[edge,dashed] (buf1) -- (memory.north -| buf1);
\draw[edge,densely dotted,gray]
  ($(core0.south east)+(-0.3,0)$) .. controls (2.6,-1.5)
  .. ($(memory.north)+(2.9,0)$);
\draw[edge,densely dotted,gray]
  ($(core1.south west)+(0.3,0)$) .. controls (4.0,-1.5)
  .. ($(memory.north)+(3.7,0)$);
\node[font=\footnotesize,gray] at (3.3,-1.35)
  {loads bypass the buffer, read memory directly};
\node[draw=red,dashed,rounded corners,align=center,font=\footnotesize,
      text width=10.6cm] (anom) at (3.3,-5.9)
 {outcome $r_1{=}0$ and $r_2{=}0$: each load bypasses its own pending store
  and reads old memory.\\ Forbidden under SC --- allowed under TSO (the
  store-buffer anomaly). A full \texttt{mfence} between the store and the
  load forbids it.};
\end{tikzpicture}
\caption{Total store order. Each core's store waits in its private store
buffer (solid arrow) and drains to shared memory later (dashed arrows);
the subsequent load bypasses the buffer (dotted arrows) and reads the old
value. The only reordering TSO permits is exactly this store-then-load
pair.}
\label{fig:tso-model}
\end{figure}

\subsection{The Relaxed Models of ARM and POWER}
\label{sec:weak-model}

ARM (before and after ARMv8) and POWER are \emph{weakly ordered}: for
ordinary loads and stores, all four pairs---store after store, load after
load, load then store, and store then load---may be reordered, subject to
per-location coherence and to data/address/control dependencies
\cite{alglave17fencing,sorin11primer}. The relaxation has several physical
sources at once: a store queue that can bypass and coalesce, write-back
caches that propagate at different times, and loads that execute
speculatively or complete out of order while waiting on misses. Ordering
is recovered only where the program asks for it: explicit barriers (ARM
\texttt{dmb}, POWER \texttt{sync}/\texttt{lwsync}/\texttt{isync}),
the ARMv8 acquire/release load-store instructions (\texttt{ldar}/
\texttt{stlr}), or the one-way ordering that dependencies already provide
(an address-dependent load cannot be hoisted above the load that computes
its address).

Figure~\ref{fig:weak-model} shows the canonical anomaly: the
\emph{message-passing} (MP) litmus test. The writer publishes a payload
and then a flag; the reader polls the flag and then reads the payload.
On ARM or POWER, with plain loads and stores, the reader can see the flag
set and still read the old payload: the two stores left the writer in a
different order, or the two loads completed out of order. The transactional
reading of the failure is exactly the commit-publish problem: the flag is
a commit marker, the payload is the write-set, and a reader that trusts
the marker without a barrier observes a torn commit. The fix is a release
store for the flag and an acquire load for it (or \texttt{dmb ish} on
ARM, \texttt{lwsync} on POWER); Section~\ref{sec:mp-example} turns this
into compilable C++ and Rust.

Weak models permit the store-buffering outcome of
Figure~\ref{fig:tso-model} as well, and one more test that TSO forbids:
IRIW (independent reads of independent writes), where two readers
disagree about the order of two independent stores. Table~\ref{tab:litmus-models}
summarizes. What survives on every platform is per-location coherence
(two writes to the same word are seen in one order by everyone) and
dependency-ordered reads---facts that the lock-free read paths of
Part~IV lean on heavily.

The repository carries one backend that targets this family directly:
a POWER8 HTM-with-fallback runtime (\texttt{backends/tm\_impl/
power8\_htm/}), which uses the hardware transactional instructions of
POWER8 and inherits exactly these ordering rules for its fallback path.
Every x86-tuned fence placement in this book was audited against the
ARM/POWER column of Table~\ref{tab:litmus-models} before the Rust and C++
runtimes shipped.

\begin{figure}[htbp]
\centering
\begin{tikzpicture}[
font=\small,
core/.style={draw,rectangle,rounded corners,align=center,
minimum width=5.0cm,minimum height=1.5cm},
mem/.style={draw,rectangle,rounded corners,align=center,
minimum width=11.2cm,minimum height=1.0cm},
edge/.style={-{Stealth},thick},
lab/.style={font=\footnotesize,fill=white,inner sep=1.5pt}
]
\node[core] (w) at (0,0)
  {Core 0 (writer)\\\texttt{W} $A{=}1$; \ then \texttt{W} $f{=}1$};
\node[core] (r) at (6.8,0)
  {Core 1 (reader)\\\texttt{R} $f \to 1$; \ then \texttt{R} $A \to 0$};
\node[mem] (memory) at (3.4,-3.5)
  {shared memory: \ $A{=}1$, $f{=}1$ \ (eventually)};
\draw[edge] (w.south) -- (memory.north -| w.south)
  node[lab, midway, right=2pt, align=left]
  {stores may leave\\in either order};
\draw[edge] (r.south) -- (memory.north -| r.south)
  node[lab, midway, left=2pt, align=right]
  {loads may complete\\out of order};
\node[draw=red,dashed,rounded corners,align=center,font=\footnotesize,
      text width=10.6cm] (anom) at (3.4,-5.2)
 {the reader sees the flag $f{=}1$ but the stale $A{=}0$: the
  ``committed'' message arrived before the data it announces.\\
  Forbidden under SC and TSO --- allowed on ARM and POWER with plain
  loads and stores. Fix: release store of $f$, acquire load of $f$
  (\texttt{stlr}/\texttt{ldar}, \texttt{dmb ish}, \texttt{lwsync}).};
\end{tikzpicture}
\caption{The weakly ordered ARM and POWER models. Both the store-store
pair on the writer and the load-load pair on the reader may be reordered,
so the message-passing idiom fails unless release/acquire ordering (or a
barrier) is inserted on at least one side.}
\label{fig:weak-model}
\end{figure}

\subsection{Other Models and the C/C++ Contract}

Two further points on the spectrum matter for later chapters. RISC-V's
RVWMO is, for practical purposes, a cousin of the ARM/POWER column of
Table~\ref{tab:litmus-models}: ordinary accesses are weakly ordered, and
acquire/release annotations plus explicit fences recover order. GPUs are
weaker still in a different dimension: the private L1 caches of a
streaming multiprocessor are not coherent across SMs at all, so
``the order in which stores become visible'' is a per-scope question;
Chapter~17 rebuilds conflict detection from scratch under that constraint.

The portable contract over all of these is the C11/C++11 atomic memory
model (Chapter~4): \texttt{seq\_cst} gives the SC column, acquire/release
gives exactly the one-way ordering needed for the message-passing idiom on
every architecture in the table, and \texttt{relaxed} gives atomicity with
no ordering at all. A compiler fence
(\texttt{atomic\_signal\_fence}) constrains only the compiler's
reordering; a hardware fence (\texttt{mfence}, \texttt{dmb},
\texttt{sync}) constrains the machine. The TM runtimes of Part~IV choose
the weakest tool at each point: compiler fences on read paths that x86
already orders, release/acquire on commit-publish paths, and full fences
only where a protocol's correctness proof requires them.

Table~\ref{tab:litmus-models} collects the litmus answers developed in
this section.

\begin{table}[htbp]
\centering
\caption{Litmus tests under the three memory-model families. Sequential
consistency forbids every surprising outcome listed, so it has no column:
``allowed'' means some execution of the litmus test may observe the
outcome. Fence and atomic fixes are discussed in the text.}
\label{tab:litmus-models}
\renewcommand{\arraystretch}{1.3}
\setlength{\tabcolsep}{4pt}
\begin{tabular}{M{0.11\linewidth} M{0.28\linewidth} M{0.23\linewidth} M{0.12\linewidth} M{0.14\linewidth}}
\hline
\thb{0.11\textwidth}{\textbf{Test}} &
\thb{0.28\textwidth}{\textbf{Program (per thread)}} &
\thb{0.23\textwidth}{\textbf{Surprising outcome}} &
\thb{0.12\textwidth}{\textbf{TSO}} &
\thb{0.14\textwidth}{\textbf{ARM / POWER}} \\
\hline
SB (store buffering) &
T1: \texttt{W} $x{=}1$; \texttt{R} $y \to r_1$.\ T2: \texttt{W} $y{=}1$; \texttt{R} $x \to r_2$. &
$r_1{=}0$ and $r_2{=}0$ (both ``went first'') &
allowed &
allowed \\
MP (message passing) &
T1: \texttt{W} $x{=}1$; \texttt{W} $f{=}1$.\ T2: \texttt{R} $f \to r_1$; \texttt{R} $x \to r_2$. &
$r_1{=}1$ but $r_2{=}0$ (flag seen, payload stale) &
forbidden &
allowed \\
IRIW (independent reads of independent writes) &
T1: \texttt{R} $x \to r_1$; \texttt{R} $y \to r_2$.\ T2: \texttt{R} $y \to r_3$; \texttt{R} $x \to r_4$. &
$r_1{=}r_3{=}1$ but $r_2{=}r_4{=}0$ (readers disagree on store order) &
forbidden &
allowed \\
\hline
\end{tabular}
\end{table}

\section{Worked Example: The TSO Store-Buffer Surprise}
\label{sec:tso-example}

The running bank example can expose a memory-ordering subtlety. Consider a
transfer that writes account $A$ and then reads account $B$:

\begin{lstlisting}[language=C, caption={A transfer that writes then reads.},
label={lst:tso-transfer-check}]
void transfer_and_check(struct account *a, struct account *b, int n) {
    a->balance -= n;            // store to A
    int b_bal = b->balance;     // load from B
    if (b_bal < 0) {
        // some recovery action
    }
    b->balance += n;            // store to B
}
\end{lstlisting}

On TSO (\S\ref{sec:tso-model}), the store to \texttt{a->balance} sits in the store
buffer. The subsequent load of \texttt{b->balance} can be satisfied before
the store to $A$ is visible to other cores. A concurrent thread can observe
$A$ updated while $B$ is still the old value---a torn view of the
transaction, shown in Figure~\ref{fig:tso-store-buffer}. The fix is a
release fence on commit: the write-back of the whole write-set must become
visible at once. This is the store-buffer echo of the snapshot-consistency
problem from Chapter~2, revisited in Part~IV for software TM and Part~VII
for GPUs.

\begin{figure}[htbp]
\centering
\begin{tikzpicture}[
font=\small,
box/.style={draw,rounded corners,align=center,
minimum width=4.6cm,minimum height=1.05cm},
edge/.style={-{Stealth},thick},
lab/.style={font=\footnotesize,fill=white,inner sep=1.5pt}
]
% left column: core 0 pipeline
\node[box] (core0) at (0,0)
  {Core 0\\transfer: \texttt{W} $A{=}90$; \ then \texttt{W} $B{=}90$};
\node[box] (buf0) at (0,-2.0)
  {store buffer\\pending: $A{=}90$, $B{=}90$};
\node[box] (memory) at (0,-4.0)
  {shared memory (draining)\\$A{=}90$, \ $B{=}100$};
% right: observing core
\node[box] (core1) at (7.0,-2.0)
  {Core 1\\reads $A$, then reads $B$};
% arrows
\draw[edge] (core0) -- (buf0)
  node[lab, midway, right=2pt] {stores enter};
\draw[edge,dashed] (buf0) -- (memory)
  node[lab, midway, right=2pt, align=left]
  {drains one store\\at a time};
\draw[edge,dashed,red] (core1.south) -- (memory.east)
  node[lab, pos=0.45, above=2pt]
  {observes the torn state};
\node[font=\footnotesize, align=left, anchor=north west]
  at ($(memory.south west)+(-0.1,-0.35)$)
  {window between the two drains: $A$ updated, $B$ not yet};
\end{tikzpicture}
\caption{TSO store-buffer surprise. Core~0 debits $A$ and credits $B$;
the two stores drain from the store buffer one at a time. A reader in the
window between the two drains observes $A{=}90$ with $B{=}100$: money has
appeared. A release fence on commit makes the whole write-set visible
atomically.}
\label{fig:tso-store-buffer}
\end{figure}

\section[The Message-Passing Surprise]{Worked Example: The Message-Passing Surprise on ARM and POWER}
\label{sec:mp-example}

The same bank transfer breaks differently on a weakly ordered machine. The
idiom is message passing (\S\ref{sec:weak-model}): the runtime publishes the write-set
and then a commit marker; readers poll the marker and then trust the
write-set. On x86 the store-store and load-load pairs are already ordered,
so the idiom survives with compiler fences alone. On ARM or POWER, plain
stores and loads permit the outcome of Figure~\ref{fig:weak-model}: a
reader sees \emph{committed} and then reads the old balance.

The listing below shows the failure and the fix in compilable C++. The
broken pair uses \texttt{relaxed} on both sides---atomicity without
ordering, which is exactly ``plain load and store'' as far as the memory
model is concerned. The fixed pair publishes the balance with a release
store and observes the marker with an acquire load, which compiles to
\texttt{stlr}/\texttt{ldar} on AArch64, \texttt{lwsync}-compatible code on
POWER, and---on x86---to nothing beyond ordinary stores and loads.

\begin{lstlisting}[language=C++, caption={Message passing with atomics:
the relaxed pair can fail on ARM/POWER; the release/acquire pair cannot.},
label={lst:mp-atomics}]
#include <atomic>

std::atomic<int> balance{100};     // the payload (the write-set)
std::atomic<bool> committed{false}; // the commit marker

void deposit_broken() {
    balance.store(90, std::memory_order_relaxed);
    committed.store(true, std::memory_order_relaxed);   // no ordering!
}

bool observe_broken(int &out) {
    if (!committed.load(std::memory_order_relaxed)) return false;
    out = balance.load(std::memory_order_relaxed);  // may read 100
    return true;
}

void deposit_fixed() {
    balance.store(90, std::memory_order_relaxed);
    committed.store(true, std::memory_order_release);   // publish
}

bool observe_fixed(int &out) {
    if (!committed.load(std::memory_order_acquire)) return false;
    out = balance.load(std::memory_order_relaxed);  // guaranteed 90
    return true;
}
\end{lstlisting}

Rust expresses the identical protocol through its atomic types, and the
type system adds one guarantee C++ cannot: a \texttt{\&T} shared across
threads cannot be mutated outside an atomic or a transaction at all, so
the ``forgot the ordering'' failure mode below is the only one left. The
repository's Rust runtimes (\texttt{explicit\_api/rust/workspace/runtime/})
use exactly these orderings: compiler fences on read paths that x86
already orders, release/acquire on the commit-publish path, and a full
fence where the backend's proof requires one.

\begin{lstlisting}[language=Rust, caption={The same message-passing
protocol in Rust: release publishes the commit, acquire observes it.},
label={lst:mp-rust}]
use std::sync::atomic::{
    AtomicI32, AtomicBool,
    Ordering::{Acquire, Relaxed, Release},
};

static BALANCE: AtomicI32 = AtomicI32::new(100); // the write-set
static COMMITTED: AtomicBool = AtomicBool::new(false); // the marker

fn deposit() {
    BALANCE.store(90, Relaxed);
    COMMITTED.store(true, Release); // publish (stlr on AArch64)
}

fn observe() -> Option<i32> {
    if !COMMITTED.load(Acquire) {   // subscribe (ldar on AArch64)
        return None;
    }
    Some(BALANCE.load(Relaxed))     // guaranteed 90, never 100
}

fn main() {
    deposit();
    assert_eq!(observe(), Some(90));
}
\end{lstlisting}

A useful discipline falls out of Table~\ref{tab:litmus-models}: write the
commit protocol against the ARM/POWER column, and treat the TSO column as
an optimization. A protocol correct on the weak column is correct
everywhere in this book's backend zoo, from x86 laptops to the POWER8
HTM backend to the GPU designs of Part~VII (where the coherence scope
shrinks further still).

\section{Worked Example: Cache-Line False Conflicts}

Hardware transactional memory normally detects conflicts at cache-line
granularity. Two independent accounts can conflict if they share a line, even
when no transaction touches the same word. The following layout check shows
the problem.

\begin{lstlisting}[language=C++, caption={Separating bank accounts onto
cache lines.}, label={lst:cache-line-layout}]
#include <atomic>
#include <cassert>
#include <cstdint>
#include <mutex>
#include <thread>
#include <vector>

constexpr int CACHE_LINE = 64;

struct Account {
    int balance;
};

struct alignas(CACHE_LINE) PaddedAccount {
    int balance;
};

static_assert(sizeof(PaddedAccount) <= CACHE_LINE,
              "PaddedAccount must not exceed one cache line");

std::mutex g_lock;

template <typename T>
void transfer_sgl(T &a, T &b, int n) {
    std::lock_guard<std::mutex> lg(g_lock);
    a.balance -= n;
    b.balance += n;
}

template <typename T>
int line_of(const T *p) {
    return (int)(reinterpret_cast<std::uintptr_t>(p) / CACHE_LINE);
}

int main() {
    Account a{100}, b{100};
    PaddedAccount pa{100}, pb{100};

    static_assert(alignof(PaddedAccount) >= CACHE_LINE,
                  "PaddedAccount must be cache-line aligned");

    if (line_of(&pa) == line_of(&pb)) {
        assert(false && "Padded accounts must live on distinct lines");
    }

    for (int i = 0; i < 1000; ++i) {
        transfer_sgl(a, b, 1);
        transfer_sgl(pa, pb, 1);
    }

    assert(a.balance + b.balance == 200);
    assert(pa.balance + pb.balance == 200);
    return 0;
}
\end{lstlisting}

Coherence works on blocks, not abstract objects. HTM conflict detection
inherits that granularity. Padding can reduce false conflicts, but it
increases footprint and can increase capacity aborts. SGL-style execution is
insensitive to false conflicts, but it also removes concurrent execution.
TSX-SGL~\cite{intel13tsx} uses HTM when the hardware accepts the transaction
and falls back to a single global lock when it does not, so it pays the
false-conflict cost only on the hardware fast path. Both live in the
companion repository: TSX-SGL in
\texttt{backends/\allowbreak tm\_impl/\allowbreak tsx\_sgl/} and
the pure-lock reference in
\texttt{backends/\allowbreak tm\_impl/\allowbreak single\_global\_lock/}.

\section{HTM Execution States}

Figure~\ref{fig:htm-state-machine} shows the state machine of a typical
best-effort hardware transactional memory implementation. Speculative data
live in the cache while the transaction is active. Coherence invalidations and
capacity overflow abort the transaction. Robust implementations include a
fallback path.

\begin{figure}[htbp]
\centering
\begin{tikzpicture}[
font=\small,
state/.style={draw,rectangle,rounded corners,align=center,
minimum width=3.2cm,minimum height=1.0cm},
edge/.style={-{Stealth},thick},
lab/.style={font=\footnotesize,fill=white,inner sep=2pt}
]
\node[state] (normal)   at (0,0)       {nontransactional\\execution};
\node[state] (active)   at (5.6,0)     {transaction active\\read/write set in cache};
\node[state] (commit)   at (12.2,0)    {commit\\publish writes};
\node[state] (abort)    at (5.6,-2.8)  {abort\\discard speculative state};
\node[state] (fallback) at (12.2,-2.8) {fallback\\SGL or STM path};

\draw[edge] (normal) -- node[lab] {xbegin} (active);
\draw[edge] (active) -- node[lab] {no conflict} (commit);
\draw[edge] (active) -- node[lab, right=3pt, align=left]
  {coherence conflict\\or capacity} (abort);
\draw[edge] (abort) -- node[lab, above left=1pt and -8pt] {retry budget} (normal);
\draw[edge] (abort) -- node[lab] {persistent abort} (fallback);
\draw[edge] (fallback.south) -- ++(0,-0.55) -|
  node[lab, pos=0.25, below] {unlock} (normal.south);
\end{tikzpicture}
\caption{HTM execution state machine. Speculative data live in the cache
while the transaction is active. Coherence invalidations and capacity
overflow abort the transaction. A robust implementation includes a
fallback path for persistent aborts.}
\label{fig:htm-state-machine}
\end{figure}

The state machine is intentionally small. The hardware does not know about
transaction bodies; it tracks cache lines and coherence messages. The
programmer, or the compiler, inserts the begin and commit instructions. The
hardware then enforces atomic visibility at commit.

\section{Worked Example: Atomic Commit-Publish in TLA+}

The hardware rule ``publish all writes at commit'' has a direct specification
counterpart: the transfer is one next-state action. The following finite TLA+
module models exactly that.

\begin{lstlisting}[caption={A finite TLA+ model of atomic bank transfers.},
label={lst:bank-atomic-commit}]
---- MODULE BankAtomicCommit ----
EXTENDS Naturals, TLC

CONSTANT Floor

Accounts == {"A", "B"}

VARIABLE balance
vars == <<balance>>

Init ==
/\ balance = [a \in Accounts |-> 100]

Transfer(a, b, n) ==
/\ a \in Accounts
/\ b \in Accounts
/\ a /= b
/\ n \in 1..10
/\ balance[a] >= Floor + n
/\ balance' = [balance EXCEPT ![a] = balance[a] - n,
![b] = balance[b] + n]

Next ==
\/ \E a,b \in Accounts, n \in 1..10 : Transfer(a,b,n)
\/ UNCHANGED vars

Total ==
balance["A"] + balance["B"] = 200

FloorOK ==
\A a \in Accounts : balance[a] >= Floor

Invariant ==
/\ Total
/\ FloorOK

Spec == Init /\ [][Next]_balance
====
\end{lstlisting}

The write-set is modeled by the function update in \texttt{balance'}.
The invariant combines money conservation and the account floor. If the
transfer is split into a visible debit action followed by a visible credit
action, TLC can reach an intermediate state that violates money conservation.
The same atomicity obligation appears in software TM write-back, HTM commit,
persistent TM log replay, and distributed two-phase commit.

\section{Hardware Mechanisms and TM Consequences}

Table~\ref{tab:hardware-tm-taxonomy} summarizes the hardware mechanisms this
chapter introduced and the TM design consequences each one forces.

\begin{table}[htbp]
\centering
\caption{Hardware mechanisms and their TM consequences. Each representative
backend lists its citation and its implementation directory in the
companion repository (Rust runtimes live under
\texttt{explicit\_api/rust/workspace/runtime/}).}
\label{tab:hardware-tm-taxonomy}
\renewcommand{\arraystretch}{1.3}
\begin{tabular}{M{0.17\linewidth} M{0.30\linewidth} M{0.41\linewidth}}
\hline
\thb{0.17\textwidth}{\textbf{Mechanism}} &
\thb{0.30\textwidth}{\textbf{TM consequence}} &
\thb{0.41\textwidth}{\textbf{Representative backends (cite, repo path)}} \\
\hline
Cache coherence &
Conflict signal on invalidation; hardware conflict detection; false conflicts at line granularity &
TSX-SGL~\cite{intel13tsx,herlihy93} (\rep{backends/\allowbreak tm\_impl/\allowbreak tsx\_sgl}); POWER8-HTM (\rep{backends/\allowbreak tm\_impl/\allowbreak power8\_htm}); all HTM paths of Ch.~16 \\
Private L1 cache &
Speculative data home; capacity aborts when transaction working set overflows &
TSX-SGL~\cite{intel13tsx}; XTM~\cite{chung06xtm} virtualizes capacity at page granularity (\rep{backends/\allowbreak tm\_impl/\allowbreak xtm}); software backends avoid it by keeping metadata in main memory (Part~IV) \\
Store buffer &
Relaxed write-to-read ordering; torn views on commit without fences; commit-publish fence required &
TL2~\cite{dice06tl2}, NOrec~\cite{dalessandro10norec}, TinySTM~\cite{felber08tinystm}, SwissTM~\cite{dragojevic09swiss}, JVSTM~\cite{cachopo06}, MVLog (companion); all under \rep{backends/\allowbreak tm\_impl/}\,$\langle$name$\rangle$ \\
Timestamp counter or logical clock &
Global commit order; version clocks for read-set validation &
TL2~\cite{dice06tl2}, TSC-TM (companion variant of TL2), NOrec-BF~\cite{dalessandro10norec} (companion extension), JVSTM~\cite{cachopo06}, MVLog (companion), LEFTRIGHT~\cite{ramalhete15leftright} (companion OCC variant, \rep{backends/\allowbreak tm\_impl/\allowbreak leftright}) \\
No coherent cache across GPU SMs &
Software conflict detection; warp-level voting; multi-version commit logs &
PR-STM~\cite{shen15prstm}, CSMV~\cite{nunes22csmv}, GPUTX~\cite{xu14gpustm}, GAccO~\cite{holey14}, GUST~\cite{nunes23gust}; \rep{gpu/\allowbreak backends/} \\
Persistent memory ordering &
Durable commit requires write ordering and crash consistency; undo/redo logs &
NV-HTM~\cite{castro18nvhtm}, Romulus~\cite{correia18romulus}, DUDETM~\cite{liu17dudetm}, PersistentSGL (companion); all under \rep{backends/\allowbreak tm\_impl/}\,$\langle$name$\rangle$ \\
Distributed memory &
No shared coherence domain; transactions become distributed protocols; two-phase commit &
DistributedSGL (companion), TiKV~\cite{peng10percolator} (\rep{backends/\allowbreak tm\_impl/\allowbreak tikv}; Rust: \rep{explicit\_api/\allowbreak rust/\allowbreak workspace/\allowbreak runtime/\allowbreak tikv}) \\
\hline
\end{tabular}
\end{table}

The table makes explicit what the rest of the book assumes: every TM backend,
whether software or hardware, is built on this substrate. The caches supply
conflict detection or force software validation. The store buffer forces the
commit-publish fence. The absence of coherent caches on GPUs forces a
different design family entirely. The presence of persistent memory adds a
durability obligation to the same atomic-publish rule.

\section{Summary}

Modern CPUs already implement a form of speculative rollback through branch
prediction and reorder buffers. Transactional memory reuses and extends that
mechanism: a transaction's speculative state lives in the cache, and abort
discards it. Cache coherence can detect conflicts because an invalidation is
a signal that a line touched by the transaction has been written elsewhere.
But store buffers hide visibility, so a software TM must insert fences to make
a commit's write-set visible atomically. The hardware substrate thus both
enables and constrains every design in the chapters that follow.

The memory-model ladder of \S3.3 gives that obligation a precise shape.
Sequential consistency has no anomalies; TSO has exactly one (the
store-buffering outcome of Figure~\ref{fig:tso-model}); the ARM and POWER
family adds message-passing and IRIW anomalies (Figure~\ref{fig:weak-model}).
A commit-publish protocol that is correct in the weakest column of
Table~\ref{tab:litmus-models} is correct everywhere in this book.

The store-buffer surprise of Section~\ref{sec:tso-example} is the hardware
echo of the snapshot-consistency problem from Chapter~2, and the
message-passing surprise of Section~\ref{sec:mp-example} is its portable
repair in both C++ and Rust. The false-conflict example shows why
coarse-grained hardware conflict detection is not free. The atomic
commit-publish TLA+ model states the obligation that every backend in
Parts~IV and VII must refine: the debit and the credit are one step, not two.

\section{Exercises}

\begin{enumerate}
\item \emph{[theory]} Draw a pipeline diagram showing a speculatively
executed store that must be squashed on an abort. Label the reorder
buffer, the memory stage, and the writeback stage. Explain why the
store never becomes architecturally visible.

\item \emph{[theory]} In a system with TSO, a transaction writes to
location A and then reads B. Show a scenario where another thread sees
the write to A before the transaction commits, violating opacity.
Propose a fence placement that prevents the torn view.

\item \emph{[theory]} Explain why a TM that uses only cache coherence for
conflict detection might still suffer from store-buffer effects.

\item \emph{[implementation]} Modify the cache-line example in
Listing~\ref{lst:cache-line-layout} so that four worker threads
repeatedly transfer between adjacent accounts. Provide two layouts:
compact and padded. Count committed transfers only under the SGL
wrapper, and check after every run that the total balance is
unchanged. Do not report timing as a research result; use the
experiment only to inspect layout and contention.

\item \emph{[verification]} Extend the TLA+ model in
Listing~\ref{lst:bank-atomic-commit} by adding a broken two-step
transfer with separate \texttt{Debit} and \texttt{Credit} actions.
State the money-conservation invariant explicitly:
$\sum_{a \in \mathrm{Accounts}} \mathrm{balance}[a] = 200$.
Then model-check that TLC finds a violation for the broken model but
not for the atomic \texttt{Transfer} action.

\item \emph{[theory]} Explain why a cache-coherence invalidation is
sufficient to detect some HTM conflicts but not sufficient to
implement a full STM on a GPU. Your answer should distinguish word
conflicts, cache-line conflicts, and the absence of coherent private
caches across GPU SMs.  (The coherence machinery itself is surveyed in
Sorin, Hill, and Wood's primer \cite{sorin11primer}.)

\item \emph{[implementation]} Write a small x86 program that reproduces
the store-buffer surprise of Section~\ref{sec:tso-example}. Use two
threads, a release fence, and an atomic flag to coordinate. Verify
that without the fence, the torn view is observable in some runs, and
with the fence, it is not.

\item \emph{[implementation]} Port the message-passing litmus test of
Listing~\ref{lst:mp-rust} to a loop that runs \texttt{deposit}/
\texttt{observe} pairs from two threads with \texttt{relaxed} ordering
on both sides, and count how often \texttt{observe} returns the stale
value. Then switch the marker to \texttt{Release}/\texttt{Acquire} and
verify the stale count drops to zero. If you have access to an AArch64
machine (for example an Apple Silicon or Raspberry Pi host), run both
variants there and on x86-64, and explain the difference using
Table~\ref{tab:litmus-models}.

\item \emph{[open]} Herlihy and Moss' 1993 proposal \cite{herlihy93} gave
HTM conflict detection ``for free'' by piggybacking on cache
coherence, yet commodity HTM arrived late, as a best-effort feature,
and remains marginal.  Argue a position: was the 1993 design ahead of
its time, or did its assumptions (L1-sized transactions, coherent
private caches) guarantee that commodity hardware could only ship a
weakened version?  Point to one concrete mechanism from this chapter
(store buffers, capacity aborts, TSO) that supports your side.
\end{enumerate}

\textbf{How to check your work.} The \emph{[implementation]} exercises are
verified with the companion: run the cache-line layouts of
Listing~\ref{lst:cache-line-layout} under \texttt{make -C benchmarks/cpp
run-bank}, reproduce the store-buffer surprise of
\S\ref{sec:tso-example}, and run the Rust litmus loop of the
message-passing exercise (the crate in
\texttt{explicit\_api/rust/workspace/} builds with \texttt{cargo build
--release}).  The \emph{[verification]} exercise is checked with
TLC against Listing~\ref{lst:bank-atomic-commit}
($\sum \mathrm{balance}[a]=200$; see Appendix~\ref{chap:pluscal-quick-ref}).
The \emph{[theory]} exercises are graded against
\S\ref{sec:tso-example} and \S\ref{sec:hardware}.  The \emph{[open]}
exercise has no single right answer; it is graded on whether you connect
your position to at least one concrete mechanism from this chapter and
address the strongest objection.

\index{data race}
\index{happens-before}
\index{compiler barrier}
\index{acquire/release}
\index{seq\_cst}
\index{snooping}
\index{directory-based coherence}
\index{invalidation}
\index{reorder buffer}
\index{branch prediction}
\index{write-back}
\index{TLA+}
\index{invariant}
\index{money conservation}
\index{single global lock}
\index{cache-line granularity}
\index{false conflict}
\index{capacity abort}
\index{commit-publish fence}
\index{store-to-load forwarding}
\index{sequential consistency}
\index{litmus test}
\index{message passing}
\index{IRIW}
\index{Dekker's example}
\index{mfence}
\index{dmb barrier}
\index{lwsync}
\index{multi-copy atomic}
\index{RVWMO}
\index{memory consistency model}


% -------------------------------------------------------------------
